% FortySecondsCV LaTeX template
% Copyright © 2019 René Wirnata <rene.wirnata@pandascience.net>
% Licensed under the 3-Clause BSD License. See LICENSE file for details.
%
% Attributions
% ------------
% * fortysecondscv is based on the twentysecondcv class by Carmine Spagnuolo 
%   (cspagnuolo@unisa.it), released under the MIT license and available under
%   https://github.com/spagnuolocarmine/TwentySecondsCurriculumVitae-LaTex
% * further attributions are indicated immediately before corresponding code

%-------------------------------------------------------------------------------
%                             ADDITIONAL PACKAGES
%-------------------------------------------------------------------------------
\documentclass[
  a4paper, 
%   showframes,
   maincolor=cvblue,
   sectioncolor=cvblue,
%  subsectioncolor=orange
%   sidebarwidth=0.4\paperwidth,
%   topbottommargin=0.03\paperheight,
%   leftrightmargin=20pt
]{fortysecondscv}

% improve word spacing and hyphenation
\usepackage{microtype}
\usepackage{ragged2e}
\usepackage{rotating}

% take care of proper font encoding
\ifxetex
	\usepackage{fontspec}
	\defaultfontfeatures{Ligatures=TeX}
% \newfontfamily\headingfont[Path = fonts/]{segoeuib.ttf} % local font
\else
	\usepackage[utf8]{inputenc}
	\usepackage[T1]{fontenc}
% \usepackage[sfdefault]{noto} % use noto google font
\fi

% enable mathematical syntax for some symbols like \varnothing
\usepackage{amssymb}

% bubble diagram configuration
\usepackage{smartdiagram}
\smartdiagramset{
  % defaut font size is \large, so adjust to harmonize with sidebar layout
  bubble center node font = \footnotesize,
  bubble node font = \footnotesize,
  % default: 4cm/2.5cm; make minimum diameter relative to sidebar size
  bubble center node size = 0.3\sidebartextwidth,
  bubble node size = 0.25\sidebartextwidth,
  distance center/other bubbles = 1.5em,
  % set center bubble color
  bubble center node color = maincolor!70,
  % define the list of colors usable in the diagram
  set color list = {maincolor!10, maincolor!40,
  maincolor!20, maincolor!60, maincolor!35},
  % sets the opacity at which the bubbles are shown
  bubble fill opacity = 0.8,
}


%-------------------------------------------------------------------------------
%                            PERSONAL INFORMATION
%-------------------------------------------------------------------------------
% profile picture
\cvprofilepic{pics/AmirMahdiAmani}
% your name
\cvname{\LARGE Amir Mahdi Amani}
% job title/career
\cvjobtitle{AI Solution Engineer\\[0.2em] Python | Computer Vision \\[0.2em] Automation}

% address in one line: prevent line breaks (may slightly overhang the sidebar)
\cvaddress{\mbox{Bernhardstr.~119, 50968~Köln, Deutschland}}
% phone number
\cvphone{+49-152-11555803}
% email address
\cvmail{amirmahdi.amani7th@gmail.com}
% personal website
\cvsite{\href{https://www.linkedin.com/in/amirmahdiamani}{LinkedIn: amirmahdiamani}}
% pgp key
\cvkey{GitHub}{github.com/am-amani}

%-------------------------------------------------------------------------------
%                              SIDEBAR 1st PAGE
%-------------------------------------------------------------------------------
% overwrite default icons and order of personal information
% \renewcommand{\personaltable}{%
% 	\begin{personal}[0.8em]
% 		\circleicon{\faKey}      & \cvkey  \\
% 		\circleicon{\faAt}       & \cvmail \\
% 		\circleicon{\faGlobe}    & \cvsite \\
% 		\circleicon{\faPhone}    & \cvphone \\
% 		\circleicon{\faEnvelope} & \cvaddress \\
% 		\circleicon{\faInfo}     & \cvbirthday \\
% 		% add another line
% 		\circleicon{\faQuestion} & \additional
% 	\end{personal}
% }

% add more profile sections to sidebar on first page
\addtofrontsidebar{
	% include gosquare national flags from https://github.com/gosquared/flags;
	% naming according to ISO 3166-1 alpha-2 country codes
	\graphicspath{{pics/flags/}}

	\profilesection{Schwerpunkte}
		\skill{\faGraduationCap}{KI-gestützte Datenverarbeitung}
		\skill{\faGraduationCap}{Computer Vision}
		\skill{\faGraduationCap}{Automatisierung}
		\skill{\faGraduationCap}{Technische Dokumentation}
		\skill{\faGraduationCap}{Systemintegration}

	\profilesection{Skills}
	\chartlabel{Programmierung:}\\
		\pointskill{}{Python}{5}
		\pointskill{}{C++}{4}
		\pointskill{}{SQL}{3}
		\pointskill{}{MATLAB / R}{3}

	\chartlabel{KI und Vision:}\\
		\pointskill{}{PyTorch / TensorFlow}{4}
		\pointskill{}{OpenCV / YOLOv5}{4}
		\pointskill{}{Bildverarbeitung}{4}

	\chartlabel{Tools und Systeme:}\\
		\pointskill{}{Linux / Docker / Git}{4}
		\pointskill{}{ROS 1/2 / Gazebo}{4}
		\pointskill{}{QGIS / Excel}{4}
}


%-------------------------------------------------------------------------------
%                              SIDEBAR 2nd PAGE
%-------------------------------------------------------------------------------
\addtobacksidebar{
	\profilesection{Kurzprofil}
	\aboutme{
KI- und Softwareingenieur mit einem M.Sc.\ in Computer Engineering mit Schwerpunkt Artificial Intelligence und Robotics. Praktische Erfahrung in Python, C++, Deep Learning, Computer Vision, Datenanalyse, technischen Tests und Dokumentation.}

	\profilesection{Sprachen}
	\barskill{}{\textbf{Deutsch} (B1)}{45}
	\barskill{}{\textbf{Englisch} (professionell)}{85}

	\profilesection{Zertifikate}
	\aboutme{
\textbf{Introduction to TensorFlow for Artificial Intelligence, Machine Learning, and Deep Learning}\\
DeepLearning.AI / Coursera, 2021\\[0.8em]
\textbf{Machine Vision Solutions with HALCON and MERLIC}\\
Stemmer Imaging Co., Niederlande, 2019}

	\profilesection{Auszeichnungen}
	\aboutme{
Erasmus+ Traineeship Scholarship, Universität Padua\\[0.6em]
Regione Veneto Scholarship, Universität Padua (2,5 Jahre)}

	\profilesection{Profile}
	\centering
	\href{https://www.linkedin.com/in/amirmahdiamani}{\includegraphics[width=0.2\sidebartextwidth]{in.png}}\quad
	\href{https://github.com/am-amani}{\includegraphics[width=0.2\sidebartextwidth]{github.png}}
}

\addtothirdsidebar{
\textbf{}\\
\begin{turn}{-90}\chartlabel{Python}\chartlabel{Computer Vision}\chartlabel{Automation}\chartlabel{Deep Learning}\chartlabel{Technical Documentation}\chartlabel{Data Analysis}\end{turn}
}

%-------------------------------------------------------------------------------
%                         TABLE ENTRIES RIGHT COLUMN
%-------------------------------------------------------------------------------
\begin{document}

\newpage
\restoregeometry
\sidebarwidth=0.35\paperwidth

\makefrontsidebar

\cvsection{Berufserfahrung}
\begin{cvtable}
	\cvitem{Okt.\ 2024 -- heute}
	{\color{cvsectioncolor}AI Solution Engineer -- Network Planning, Optimization \& Automation Specialist}
	{\textbf{Deutsche Netzplanung GmbH, Köln}}
	{Unterstützung datenbasierter Planungs- und Optimierungsprojekte für Telekommunikationsinfrastruktur. Analyse von Flächen- und technischen Daten mit QGIS zur Klassifizierung von Regionen für den Glasfaserausbau. Erstellung technischer Berichte und Präsentationen sowie Zusammenarbeit mit Ingenieur:innen bei der Entwicklung und Umsetzung technischer Lösungen.}

	\cvitem{März 2024 -- Sep.\ 2024}
	{\color{cvsectioncolor}Fiber Network Planning \& Optimization (Teilzeit)}
	{\textbf{Deutsche Netzplanung GmbH, Köln}}
	{Unterstützung bei der Analyse technischer Daten und der Planung von Glasfaserinfrastruktur. Nutzung von QGIS zur Flächenanalyse und zur Unterstützung effizienter Netzrollout-Entscheidungen.}

	\cvitem{Feb.\ 2024 -- Sep.\ 2024}
	{\color{cvsectioncolor}Wissenschaftliche Hilfskraft -- Mobile Robot Systems}
	{\textbf{Humanoid Robots Lab, Universität Bonn}}
	{Installation und Konfiguration von ROS-Paketen, darunter SLAM und AMCL, auf einem Unitree GO1. Tests verschiedener Betriebsmodi sowie Erstellung technischer Nutzungs- und Prozessdokumentationen. Modernisierung von TurtleBot2-Systemen für Lokalisierung, Navigation und Mapping.}

	\cvitem{Sep.\ 2019 -- Nov.\ 2019}
	{\color{cvsectioncolor}Research Internship -- Design-to-Robotic-Production and Operation}
	{\textbf{TU Delft, Niederlande}}
	{Entwicklung ROS-basierter Lösungen für Forschungsprojekte. Arbeit mit MoveIt, Kinect V2, Robotik-Kinematik sowie RGB-D-basierter Objektanalyse mit Kalibrierung, Datenfilterung, Segmentierung und Clustering.}
\end{cvtable}

\cvsection{Zielrelevantes Profil}
\begin{cvtable}
	\cvitem{}
	{\color{cvsectioncolor}Übertragbare Kompetenzen für Input Management und KI-gestützte Dokumentenverarbeitung}
	{}
	{Erfahrung in der Entwicklung und dem Testen technischer Lösungen mit Python und C++, in der Verarbeitung visueller und technischer Daten, im Modelltraining, in Fehleranalyse und in technischer Dokumentation. Ziel ist es, diese Kompetenzen gezielt in OCR-, Klassifikations-, Extraktions- und Automatisierungsprozesse einzubringen und weiterzuentwickeln.}
\end{cvtable}

\cvsection{Ausbildung}
\begin{cvtable}
	\cvitem{Sep.\ 2021 -- Sep.\ 2024}
	{\color{cvsectioncolor}M.Sc.\ Computer Engineering -- Schwerpunkt Artificial Intelligence \& Robotics}
	{\textbf{Universität Padua, Italien}}
	{Abschlussnote: 96/110. Relevante Inhalte: Machine Learning, Deep Learning, Computer Vision, intelligente und industrielle Robotik sowie 3D-Datenverarbeitung.}

	\cvitem{Dez.\ 2023 -- Sep.\ 2024}
	{\color{cvsectioncolor}Erasmus+ Forschungs- und Thesisaufenthalt}
	{\textbf{Universität Bonn, Deutschland}}
	{Masterarbeit und Forschungspraktikum am Humanoid Robots Lab.}

	\cvitem{Sep.\ 2015 -- Juni 2020}
	{\color{cvsectioncolor}B.Sc.\ Software Engineering}
	{\textbf{Universität Birjand, Iran}}
	{Relevante Inhalte: C++, Datenstrukturen, Algorithmen, Robotik und Bildverarbeitung.}
\end{cvtable}

\newpage
\makebacksidebar

\cvsection{Forschung und technische Projekte}
\begin{cvtable}
	\cvitem{Dez.\ 2023 -- Sep.\ 2024}
	{\color{cvsectioncolor}Masterarbeit: Multi-Sensor Robotic Navigation Using Teacher-Student Framework}
	{\textbf{Universität Bonn}}
	{Entwicklung eines Python-basierten Navigationssystems, das wahlweise Bild- oder LiDAR-Daten verarbeitet. Implementierung eines Teacher-Student-Ansatzes mit Deep Reinforcement Learning (TD3) und überwachtem Lernen. Umfangreiche Tests in ROS/Gazebo auf Pioneer 3-DX und TurtleBot3 Waffle sowie Auswertung der Robustheit unter unterschiedlichen Sensormodalitäten.}

	\cvitem{Juni 2022}
	{\color{cvsectioncolor}Human Hands Detection and Segmentation}
	{\textbf{Universitätsprojekt}}
	{Training eines YOLOv5-basierten Deep-Learning-Modells in PyTorch zur Handerkennung mit berichteter Genauigkeit von bis zu 95\%. Ergänzende C++-Segmentierung mit Konturerkennung, Thresholding und morphologischen Operationen.}

	\cvitem{Dez.\ 2023}
	{\color{cvsectioncolor}Real-Time Camera Pose and Tumor Detection System for Endoscopic Procedures}
	{\textbf{Universitätsprojekt}}
	{Entwicklung eines Echtzeit-Systems zur Kamerabewegungsbestimmung und Erkennung potenzieller Tumorregionen mit HSV-Filterung, Canny-Kantendetektion, Hough Circle Transform, Background Subtraction und morphologischen Operationen.}

	\cvitem{Apr.\ 2023}
	{\color{cvsectioncolor}Semi Global Matching for Stereo Vision}
	{\textbf{Universitätsprojekt}}
	{Implementierung von Semi Global Matching in C++ mit der OpenCV-Stereo-Bibliothek zur Berechnung von Tiefenkarten. Arbeit mit Census Transform, Hamming-Distanzen und kostenbasierter Optimierung.}
\end{cvtable}

\cvsection{Publikation}
\begin{cvtable}
	\cvitem{2026}
	{\color{cvsectioncolor}A Unified Conditional Policy for Multi-Robot Navigation via LiDAR-to-Vision Distillation}
	{\textbf{Machines (MDPI)}}
	{Amani, A.\ M., Amani, S.\ \& MajidiRad, A.\ H. Manuskript nach Peer Review in Überarbeitung.}
\end{cvtable}

\cvsection{Weitere Qualifikationen}
\begin{cvtable}
	\cvitem{März 2023}
	{\color{cvsectioncolor}Speaker -- 15th International Petroleum Technology Conference}
	{\textbf{Bangkok, Thailand}}
	{Präsentation eines wissenschaftlichen Beitrags im Namen eines Peers; Vermittlung komplexer technischer Inhalte an ein internationales Publikum.}

	\cvitem{}
	{\color{cvsectioncolor}Teaching Experience}
	{\textbf{Universität Birjand}}
	{Tutor in den Kursen Differential Equations und General Mathematics I.}

	\cvitem{}
	{\color{cvsectioncolor}Referenzen}
	{\textbf{Auf Anfrage}}
	{Prof.\ AmirHossein MajidiRad, Prof.\ Emanuele Menegatti, Prof.\ Maren Bennewitz, Prof.\ Alberto Pretto und Prof.\ Nasim Nasrabadi.}
\end{cvtable}

\end{document}
